\section{Related Work}
\label{sec:related_work}

\subsection{Long-Horizon Visuomotor Imitation Learning}

Long-horizon manipulation requires coherent behavior over extended execution.
Existing methods shorten the control horizon through action chunking,
autoregressive prediction, or diffusion-based
action generation~\citep{Zhao2023Learning,chi2024diffusionpolicy,
Haldar2025BAKU,reuss2024multimodal}. Other approaches introduce higher-level
task structure through subgoals, skill composition, progress estimation, or
planning~\citep{shiarlis2018taco,gupta2019relay,Chen2023Sequential,
Kang2025Incorporating,Jain2025LongerHorizon,Liu2026PALM,
Zhang2026AtomicVLA,Feng2025ReflectivePlanning,Zhou2026Act2Goal}.
These methods organize actions or task stages over time, but generally do not
supervise how task-relevant objects evolve from historical to target states.

\subsection{Temporal Modeling in Robot Policies}

Temporal context is commonly incorporated through observation histories,
latent memories, recurrent features, or retrieved past
frames~\citep{Liang2021Finite,Fang2025SAM2Act,shi2025memoryvla,
Mark2026BPP}. Recent policies further maintain recurrent or hierarchical
memory to capture dependencies beyond the current
observation~\citep{Ji2026HiMe,Xiao2026AVAVLA,Li2026OptimusVLA}.

A complementary direction uses future information as guidance.
Existing methods condition policies on goal observations or latent plans
~\citep{lynch2020play,Mandlekar2020Learning}, predict future visual states or
joint video-action trajectories~\citep{Hu2025VPP,Zhu2025UWM}, or learn
predictive representations from future transitions and interaction targets
~\citep{Zeng2024MPI}. Despite using historical and future context, these
approaches represent time through holistic observations, latent memories,
predicted frames, or task-level targets rather than identity-consistent state
transitions of individual objects.

\subsection{Object-Centric Representations for Manipulation}

Object-centric policies organize perception and action around manipulable
entities. Existing approaches use object
proposals or 3D object features as policy inputs
~\citep{zhu2022viola,zhu2023learning}, learn affordance-, flow-, or
contact-based interaction representations
~\citep{belkhale2023plato,xu2024im2flow2act,Zhang2026AffordGen,
Lee2025ViTaSCOPE}, or represent manipulation through explicit object motion,
including SE(3) pose trajectories and object pointflow
~\citep{Hsu2025SPOT,Yang2025PPI}. Recent work also introduces structured
object-level representations based on physical concepts or
semantic operations~\citep{Wei2026PhysicalConcepts,Tang2026SEAM}.

Object-centric world models use entity-level representations for dynamics and
future prediction
~\citep{ferraro2025focus,jeong2025ocwm}, while graph-based
approaches preserve explicit objects and relations for manipulation planning
~\citep{zhu2026orion}. These methods typically represent objects as perceptual
units, interaction interfaces, motion variables, or future prediction targets.
In contrast, OST-Map maintains identity-consistent structural states across
observed and future task stages, making past-to-future object-state transitions
explicit supervision for action generation.
